\section*{Capítulo \ref{ch:g2:water-on-earth} --- A água na Terra}

\begin{solution}{exo:g2:water-on-earth:1}
\omterm{def:g2:water-on-earth:salt-fresh}{Água salgada}: o \omterm{def:g2:water-on-earth:ocean}{mar}, o \omterm{def:g2:water-on-earth:ocean}{oceano}. \omterm{def:g2:water-on-earth:salt-fresh}{Água doce}: o \omterm{def:g1:landscapes-around-us:water}{rio}, a chuva, o \omterm{def:g1:landscapes-around-us:water}{lago} nas
\omterm{def:g1:landscapes-around-us:landform}{montanhas}, o gelo da \omterm{def:g1:landscapes-around-us:landform}{montanha}.
\end{solution}

\begin{solution}{exo:g2:water-on-earth:2}
Porque a maior parte da sua superfície é coberta pela água azul dos
\omterm{def:g2:water-on-earth:ocean}{oceanos}: vista do espaço, ela parece azul.
\end{solution}

\begin{solution}{exo:g2:water-on-earth:3}
$100 - 71 = 29$ quadrados são de terra.
\end{solution}

\begin{solution}{exo:g2:water-on-earth:4}
Um \omterm{def:g2:water-on-earth:ocean}{oceano} é maior: um \omterm{def:g2:water-on-earth:ocean}{mar} é uma extensão menor de \omterm{def:g2:water-on-earth:salt-fresh}{água salgada}, muitas
vezes na beira de um \omterm{def:g2:water-on-earth:ocean}{oceano}.
\end{solution}

\begin{solution}{exo:g2:water-on-earth:5}
Não: a água do \omterm{def:g2:water-on-earth:ocean}{oceano} é salgada, e bebê-la faz o corpo perder água. O
marinheiro precisa levar \omterm{def:g2:water-on-earth:salt-fresh}{água doce}.
\end{solution}

\begin{solution}{exo:g2:water-on-earth:6}
Por exemplo: nos \omterm{def:g1:landscapes-around-us:water}{rios}, nos \omterm{def:g1:landscapes-around-us:water}{lagos}, no gelo e na neve das \omterm{def:g1:landscapes-around-us:landform}{montanhas} e da
Antártida, embaixo da terra, na chuva.
\end{solution}

\begin{solution}{exo:g2:water-on-earth:7}
\omterm{def:g2:water-on-earth:salt-fresh}{Água salgada}: 96 baldes. \omterm{def:g2:water-on-earth:salt-fresh}{Água doce}: $2 + 2 = 4$ baldes.
(A rigor, uma parte da água embaixo da terra também é salgada; o desenho
a conta junto com a \omterm{def:g2:water-on-earth:salt-fresh}{água doce}.)
\end{solution}

\begin{solution}{exo:g2:water-on-earth:8}
Os \omterm{def:g2:water-on-earth:ocean}{oceanos} (96 baldes), depois o gelo e a neve (2 baldes), depois os \omterm{def:g1:landscapes-around-us:water}{rios}
(muito menos que um balde).
\end{solution}

\begin{solution}{exo:g2:water-on-earth:9}
Ela pode entrar no chão; pode escorrer pelo chão até um riacho ou um \omterm{def:g1:landscapes-around-us:water}{rio};
pode evaporar e voltar para o ar quando o sol sai.
\end{solution}

\begin{solution}{exo:g2:water-on-earth:10}
Sobra sal no fundo da panela. A água evaporou: ela foi para o ar em forma
de vapor de água.
\end{solution}

\begin{solution}{exo:g2:water-on-earth:11}
Por exemplo: a gota escorre pelo telhado até a calha e o bueiro, depois
vai para um riacho e um \omterm{def:g1:landscapes-around-us:water}{rio}; o \omterm{def:g1:landscapes-around-us:water}{rio} a leva até o \omterm{def:g2:water-on-earth:ocean}{mar}; num dia de sol, ela
evapora do \omterm{def:g2:water-on-earth:ocean}{mar}, sobe para o ar e passa a fazer parte de uma nuvem.
\end{solution}
